%----------------------------------------------------------------------------------------
% SOMFORMULES. Alle opgaven zijn eigen werk.
%
% Opbouw: Start, leerstof (kern, uitgewerkt voorbeeld, korte oefening), en daarna op een
% nieuwe pagina de Opdrachten (door elkaar, stijgende moeilijkheid), Koppel en Test jezelf.
% Labels van dit hoofdstuk beginnen met h3-.
%----------------------------------------------------------------------------------------

\mychapter{3}{Somformules}

\lkr{Leerplan B+S: LPD 28, formule voor de som van de eerste $n$ termen, met samengestelde intrest als context. Twee bewijzen die leerlingen moeten kunnen volgen en navertellen: de som tweemaal schrijven, en $q \cdot s_n - s_n$. De rode draad in de opdrachten: eerst herkennen welke soort rij het is, dan het aantal termen bepalen, dan pas de formule gebruiken. Oneindige sommen horen niet bij dit leerplan.}

Hoeveel heb je gespaard als je elke week wat meer opzijlegt? Hoeveel stoelen staan er in een zaal waar elke rij twee stoelen meer telt dan de vorige? Telkens tel je de termen van een rij op. Met tien termen kan dat nog met de hand, maar met honderd of duizend termen niet meer. Voor rekenkundige en meetkundige rijen bestaat er gelukkig een formule voor die som. In dit hoofdstuk leer je die formules, en vooral waarom ze kloppen: elk berust op een slimme truc die je zelf kunt navertellen.

\begin{oefeningen}[Start]
\oef{h3-start-1}

\oef{h3-verken}
\end{oefeningen}

\lkr{Laat oefening \ref{oef:h3-verken} maken vóór de uitleg. Vraag waarom je deelt door 2: dat is de stap die leerlingen het vaakst vergeten als ze de formule later moeten verklaren.}

%========================================================================================
\section{Som van een rekenkundige rij}

\subsection{De somformule}

In de verkenning schreef je de bedragen twee keer onder elkaar, de tweede keer in omgekeerde volgorde. Elke kolom gaf dan hetzelfde getal, en het gespaarde bedrag was de helft van alle kolommen samen. Die truc werkt bij elke rekenkundige rij. Eerst geven we de som een naam.

\begin{theorie}
De \emph{som van de eerste $n$ termen} van een rij is

\qquad $\formule{s_n = u_1 + u_2 + \cdots + u_n}$

\begin{stelling}{Som van een rekenkundige rij}
\begin{minipage}[t]{.6\linewidth}
\vspace{0pt}
De som van de eerste $n$ termen van een rekenkundige rij is $n$ keer het gemiddelde van de eerste en de laatste term:

\vspace{1mm}
\qquad $\formule{s_n = n \cdot \dfrac{u_1 + u_n}{2}}$
\end{minipage}\hfill
\begin{minipage}[t]{.34\linewidth}
\vspace{0pt}
\begin{tikzpicture}[x=5mm,y=2.5mm,font=\scriptsize]
  \foreach \i/\h in {0/2,1/3,2/4,3/5,4/6}{
    \fill[accent!45] (\i,0) rectangle (\i+1,\h);
    \draw[accent] (\i,0) rectangle (\i+1,\h);
    \draw[accent] (\i,\h) rectangle (\i+1,8);}
  \draw[decorate,decoration={brace,mirror,amplitude=3pt},grijs] (0,-.3) -- (5,-.3) node[midway,below=2pt,text=black] {$n$ kolommen};
  \draw[decorate,decoration={brace,mirror,amplitude=3pt},grijs] (5.15,0) -- (5.15,8) node[midway,right=2pt,text=black] {$u_1 + u_n$};
\end{tikzpicture}
\end{minipage}
\end{stelling}

\begin{bewijs}
Schrijf de som twee keer, de tweede keer in omgekeerde volgorde, en tel kolom per kolom op:
\bewijsregel{\begin{array}{r@{\;}c@{\;}c@{\;}c@{\;}c@{\;}c@{\;}c@{\;}c}
s_n &=& u_1 &+& u_2 &+ \cdots +& u_n \\
s_n &=& u_n &+& u_{n-1} &+ \cdots +& u_1 \\ \hline
2s_n &=& (u_1 + u_n) &+& (u_2 + u_{n-1}) &+ \cdots +& (u_n + u_1)
\end{array}}
Van kolom tot kolom komt er bovenaan $v$ bij en gaat er onderaan $v$ af, dus elke kolom is gelijk aan $u_1 + u_n$. Er zijn $n$ kolommen, zodat
\bewijsregel{2s_n = n\,(u_1 + u_n) \quad\Rightarrow\quad s_n = n \cdot \dfrac{u_1 + u_n}{2} .}
\end{bewijs}
\end{theorie}

De figuur toont hetzelfde: twee keer dezelfde trap vormt een rechthoek van $n$ op $u_1 + u_n$, en de som is de helft van die rechthoek. Zo kun je de formule ook lezen als: aantal termen maal gemiddelde term.

\voorbeeld{}{Bereken $s_{20}$ van de rij $7,\ 11,\ 15,\ \ldots$\label{vb:h3-rr}}{%
\uitleg{De formule vraagt de eerste en de laatste term. De laatste term, $u_{20}$, staat er niet: die berekenen we eerst met het expliciet voorschrift.}
\begin{stappen}
v = 11 - 7 = 4 & De rij is rekenkundig, met verschil 4. Controle: $15 - 11 = 4$. \\
u_{20} = 7 + 19 \cdot 4 = 83 & Van $u_1$ naar $u_{20}$ zijn het 19 sprongen. \\[1ex]
s_{20} = 20 \cdot \dfrac{7 + 83}{2} = 900 & De gemiddelde term is $\tfrac{7 + 83}{2} = 45$, en er zijn 20 termen.
\end{stappen}
\naast{De figuur toont de twintig termen als staven. De stippellijn ligt op de gemiddelde term 45: wat de staven rechts erboven uitsteken, vult precies de ruimte links eronder. De som is dus even groot als twintig staven van 45.}{\begin{tikzpicture}[x=.26cm,y=.022cm,font=\scriptsize]
  \foreach \i in {1,...,20}{\pgfmathsetmacro{\h}{3+4*\i}\draw[accent,fill=accent!22] (\i-1,0) rectangle (\i,\h);}
  \draw[grijs,dashed,thick] (0,45) -- (20,45) node[right,text=black] {45};
  \draw[->,grijs] (0,0) -- (20.8,0) node[right,text=black] {$n$};
  \node[below,inner sep=1.5pt] at (.5,0) {1};
  \node[below,inner sep=1.5pt] at (19.5,0) {20};
  \node[above,inner sep=1pt] at (.5,7) {7};
  \node[above,inner sep=1pt] at (19.5,83) {83};
\end{tikzpicture}}}

\nujij{Bereken $s_{30}$ van de rij $20,\ 18,\ 16,\ \ldots$}{$v = -2$, $u_{30} = -38$, $s_{30} = 30 \cdot \tfrac{20 - 38}{2} = -270$.}

\thuis{Vul je $u_n = u_1 + (n-1)v$ in, dan krijg je een tweede vorm: $s_n = \tfrac{n}{2}\bigl(2u_1 + (n-1)v\bigr)$. Die is handig als je $u_n$ nog niet kent.}

\begin{oefeningen}
\oef{h3-k1}
\end{oefeningen}

\subsection{Eerst het aantal termen}

Vaak staat een som er met een laatste term, maar zonder het aantal termen. De formule heeft dat aantal wel nodig, dus dan moet je het eerst zelf zoeken.

\voorbeeld{}{Bereken de som van alle veelvouden van 6 tussen 100 en 400.\label{vb:h3-aantal}}{%
\uitleg{De veelvouden van 6 vormen een rekenkundige rij met $v = 6$. We zoeken eerst de eerste en de laatste term, dan het aantal termen, en pas daarna de som.}
\begin{stappen}
102 + 108 + \cdots + 396 & Het eerste veelvoud boven 100 is $102 = 6 \cdot 17$, het laatste onder 400 is $396 = 6 \cdot 66$. \\
102 + (n-1) \cdot 6 = 396 & De laatste term is $u_n$. We lossen op naar $n$. \\
n - 1 = 49 \ \Rightarrow \ n = 50 & Er zijn 49 sprongen, dus 50 termen. Ook zo: van $6 \cdot 17$ tot $6 \cdot 66$ zijn het $66 - 17 + 1 = 50$ veelvouden. \\[1ex]
s_{50} = 50 \cdot \dfrac{102 + 396}{2} = 12\,450 & Nu pas gebruiken we de somformule.
\end{stappen}}

\begin{werkwijze}
Is de laatste term gegeven en niet het aantal termen, bereken dan eerst $n$ met het expliciet voorschrift. Het aantal termen is één meer dan het aantal sprongen.
\end{werkwijze}

\nujij{Bereken $50 + 54 + 58 + \cdots + 250$.}{$n = 51$, en de som is $51 \cdot 150 = 7650$.}

\begin{oefeningen}
\oef{h3-k2}
\end{oefeningen}

\subsection{Het aantal termen uit de som}

Je kunt de vraag ook omdraaien: de som is gegeven, en je zoekt hoeveel termen je moet optellen. Dan is $n$ de onbekende in de somformule.

\voorbeeld{}{Hoeveel termen van $3,\ 7,\ 11,\ \ldots$ moet je optellen om 210 te krijgen?\label{vb:h3-nuit}}{%
\uitleg{In de somformule staat $u_n$, en die hangt ook af van $n$. We schrijven dus eerst $u_n$ met $n$, en daarna de hele som. Omdat $n$ dan twee keer voorkomt, krijgen we een vierkantsvergelijking.}
\begin{stappen}
u_n = 3 + (n-1) \cdot 4 = 4n - 1 & De laatste term, met $n$ als onbekende. \\[1ex]
s_n = n \cdot \dfrac{3 + 4n - 1}{2} = n\,(2n + 1) & We vullen $u_1$ en $u_n$ in de somformule in. \\
2n^2 + n - 210 = 0 & We stellen $s_n$ gelijk aan 210. \\[1ex]
n = 10 \quad\text{of}\quad n = -\dfrac{21}{2} & Alleen $n = 10$ is een aantal termen. Controle: $10 \cdot 21 = 210$.
\end{stappen}}

\begin{werkwijze}
Is de som gegeven, schrijf dan $s_n$ met $n$ als onbekende. Dat geeft een vierkantsvergelijking. Alleen een oplossing in $\mathbb{N}_0$ is een aantal termen.
\end{werkwijze}

\nujij{Hoeveel termen van $6,\ 8,\ 10,\ \ldots$ moet je optellen om 150 te krijgen?}{$s_n = n(n+5)$, dus $n = 10$.}

\begin{oefeningen}
\oef{h3-k3}
\end{oefeningen}

%========================================================================================
\section{Som van een meetkundige rij}

Bij een meetkundige rij helpt omkeren niet: de kolommen zijn dan niet meer gelijk. Er is wel een andere truc. Probeer ze eerst zelf.

\begin{oefeningen}
\oef{h3-mverken}
\end{oefeningen}

\lkr{Laat de leerlingen vergelijken met de spaaroefening: bij een rekenkundige rij kies je voor optellen in omgekeerde volgorde, bij een meetkundige rij voor vermenigvuldigen met $q$. Vraag waarom de truc van de ene niet werkt bij de andere.}

\subsection{De somformule}

In de verkenning vermenigvuldigde je de som met 3. Daardoor schoof elke term een plaats op, en vielen bij het aftrekken bijna alle termen weg. Hetzelfde lukt met elk quotiënt $q$.

\begin{theorie}
\begin{stelling}{Som van een meetkundige rij}
Voor een meetkundige rij met $q \neq 1$ is

\vspace{1mm}
\qquad $\formule{s_n = u_1 \cdot \dfrac{q^{\,n} - 1}{q - 1}}$
\end{stelling}

\begin{bewijs}
Schrijf de som, en daaronder de som maal $q$. Elke term schuift dan een plaats op. Trek de bovenste van de onderste af: alle termen in het midden vallen weg.
\bewijsregel{\begin{array}{r@{\;}c@{\;}c@{\;}c@{\;}c@{\;}c@{\;}c@{\;}c}
s_n &=& u_1 &+& u_1 q &+ \cdots +& u_1 q^{\,n-1} & \\
q\,s_n &=& &&  u_1 q &+ \cdots +& u_1 q^{\,n-1} &+\ u_1 q^{\,n} \\ \hline
q\,s_n - s_n &=& -u_1 &&&&& +\ u_1 q^{\,n}
\end{array}}
Brengen we links $s_n$ en rechts $u_1$ buiten haakjes, dan is
\bewijsregel{s_n\,(q - 1) = u_1\,(q^{\,n} - 1) .}
Omdat $q \neq 1$, mogen we delen door $q - 1$, en dat geeft de formule.
\end{bewijs}
\end{theorie}

Voor $q = 1$ zijn alle termen gelijk aan $u_1$, en is de som gewoon $n \cdot u_1$. Let ook op de exponent: in de somformule staat $n$, het aantal termen. In het expliciet voorschrift staat $n - 1$, het aantal sprongen.

\voorbeeld{}{Bereken $s_{10}$ van de rij $3,\ 6,\ 12,\ \ldots$\label{vb:h3-mr}}{%
\uitleg{Elke term is het dubbele van de vorige, dus we gebruiken de somformule voor een meetkundige rij. Daarvoor hebben we alleen $u_1$, $q$ en het aantal termen nodig, en niet de laatste term.}
\begin{stappen}
q = \dfrac{6}{3} = 2 & De rij is meetkundig, met quotiënt 2. Controle: $\dfrac{12}{6} = 2$. \\[1.5ex]
s_{10} = 3 \cdot \dfrac{2^{10} - 1}{2 - 1} = 3069 & De exponent is 10, want we tellen 10 termen op.
\end{stappen}}

\nujij{Bereken $s_9$ van de rij $1,\ {-2},\ 4,\ \ldots$}{$q = -2$, dus $s_9 = \tfrac{(-2)^9 - 1}{-2 - 1} = \tfrac{-513}{-3} = 171$.}

\begin{oefeningen}
\oef{h3-k4}
\end{oefeningen}

\subsection{Eerst het aantal termen}

Ook bij een meetkundige rij staat er soms een laatste term in plaats van een aantal termen. Dan zoek je het aantal eerst met het expliciet voorschrift, net zoals bij een rekenkundige rij. Alleen staat $n$ nu in een exponent.

\voorbeeld{}{Bereken $2 + 6 + 18 + \cdots + 4374$.\label{vb:h3-mrn}}{%
\uitleg{De rij is meetkundig met $q = 3$. We kennen de laatste term, maar niet het aantal termen. Dat zoeken we eerst, en pas dan gebruiken we de somformule.}
\begin{stappen}
2 \cdot 3^{\,n-1} = 4374 & De laatste term is $u_n = 2 \cdot 3^{\,n-1}$. \\
3^{\,n-1} = 2187 & We delen beide leden door 2. \\
n - 1 = 7 \ \Rightarrow \ n = 8 & Omdat $3^7 = 2187$, zijn het 7 sprongen en dus 8 termen. \\[1ex]
s_8 = 2 \cdot \dfrac{3^8 - 1}{3 - 1} = 6560 & De somformule met 8 termen.
\end{stappen}}

\begin{werkwijze}
Is de laatste term gegeven, bereken dan eerst $n$ uit $u_n = u_1 \cdot q^{\,n-1}$, zo nodig met een logaritme. Werk met de exacte machten, niet met afgeronde getallen.
\end{werkwijze}

\nujij{Bereken $7 + 14 + 28 + \cdots + 448$.}{$2^{n-1} = 64$, dus $n = 7$, en de som is $7 \cdot (2^7 - 1) = 889$.}

\begin{oefeningen}
\oef{h3-k5}
\end{oefeningen}

%========================================================================================
\section{Het somteken}

\subsection{Lezen en berekenen}

Een som met veel termen voluit schrijven, is lang en onoverzichtelijk. Wiskundigen gebruiken daarom een kort teken: de Griekse hoofdletter sigma, de S van \textit{som}.

\begin{theorie}
Met het \emph{somteken} $\Sigma$ schrijf je een som kort:

\qquad $\formule{\displaystyle\sum_{i=1}^{n} u_i = u_1 + u_2 + \cdots + u_n}$

\vspace{1mm}
De letter $i$ is de \emph{lopende index}: die neemt elk geheel getal aan van de beginwaarde tot en met de eindwaarde.

\begin{stelling}{Rekenregels voor het somteken}
\qquad $\formule{\displaystyle\sum_{i=1}^{n} (u_i + w_i) = \sum_{i=1}^{n} u_i + \sum_{i=1}^{n} w_i}$ \qquad\qquad $\formule{\displaystyle\sum_{i=1}^{n} c \cdot u_i = c \cdot \sum_{i=1}^{n} u_i}$
\end{stelling}

\begin{bewijs}
Schrijf beide leden voluit. Links staat $(u_1 + w_1) + \cdots + (u_n + w_n)$: bij een som mag je de termen in elke volgorde optellen, dus dat is $(u_1 + \cdots + u_n) + (w_1 + \cdots + w_n)$. In $c\,u_1 + \cdots + c\,u_n$ breng je $c$ buiten haakjes.
\end{bewijs}
\end{theorie}

Je leest $\sum_{i=1}^{n} u_i$ als: de som van de $u_i$ voor $i$ van 1 tot $n$. De naam van de index maakt niet uit: $\sum_{i=1}^{5} i$ en $\sum_{k=1}^{5} k$ zijn hetzelfde getal. De index hoeft ook niet bij 1 te beginnen. Van $i = 3$ tot $i = 12$ zijn er $12 - 3 + 1 = 10$ termen: niet 12 en niet 9.

\voorbeeld{}{Bereken $\displaystyle\sum_{i=3}^{12} (2i + 1)$.\label{vb:h3-somteken}}{%
\uitleg{We weten nog niet om welke soort rij het gaat. Daarom schrijven we eerst een paar termen uit, door de waarden van $i$ in te vullen.}
\begin{stappen}
7 + 9 + 11 + \cdots + 25 & Voor $i = 3$, 4 en 5 krijgen we 7, 9 en 11, en voor $i = 12$ krijgen we 25. Er komt telkens 2 bij: de rij is rekenkundig. \\
12 - 3 + 1 = 10 \text{ termen} & De index loopt van 3 tot en met 12. \\[1ex]
10 \cdot \dfrac{7 + 25}{2} = 160 & De somformule voor een rekenkundige rij.
\end{stappen}}

\voorbeeld{}{Bereken $\displaystyle\sum_{j=1}^{5} 3 \cdot 2^{\,j}$.\label{vb:h3-somteken2}}{%
\uitleg{Opnieuw schrijven we eerst termen uit. Let op de eerste term: die is niet 3, want de index begint bij $j = 1$, en dan staat er al $2^1$.}
\begin{stappen}
6 + 12 + 24 + 48 + 96 & Voor $j = 1$, 2 en 3 krijgen we 6, 12 en 24. Elke term is het dubbele van de vorige: de rij is meetkundig, met $u_1 = 6$ en $q = 2$. \\[1ex]
6 \cdot \dfrac{2^5 - 1}{2 - 1} = 186 & De somformule met 5 termen.
\end{stappen}}

\begin{werkwijze}
Schrijf de eerste termen uit door de index in te vullen. Herken de soort rij, tel het aantal termen (eindwaarde $-$ beginwaarde $+ 1$), en gebruik dan de somformule.
\end{werkwijze}

\lkr{Laat leerlingen altijd de eerste drie termen uitschrijven voor ze een formule kiezen. Dat voorkomt de twee klassieke fouten: het verkeerde aantal termen, en $u_1$ aflezen zonder invullen.}

\nujij{Bereken $\displaystyle\sum_{k=2}^{9} (5k - 3)$.}{Termen $7,\ 12,\ \ldots,\ 42$; 8 termen; de som is $8 \cdot \tfrac{7 + 42}{2} = 196$.}

\begin{oefeningen}
\oef{h3-k6}
\end{oefeningen}

\subsection{De som als rij}

De sommen $s_1,\ s_2,\ s_3,\ \ldots$ vormen zelf een nieuwe rij. Ken je die rij van sommen, dan kun je de oorspronkelijke termen terugvinden: het verschil tussen twee opeenvolgende sommen is precies de term die erbij kwam.

\begin{theorie}
\begin{stelling}{Termen uit de sommen}
\qquad $\formule{u_1 = s_1} \quad\text{en}\quad \formule{u_n = s_n - s_{n-1}} \quad\text{voor } n \geqslant 2$
\end{stelling}

\begin{bewijs}
De som $s_n$ is $s_{n-1}$ met de term $u_n$ erbij: $s_n = (u_1 + \cdots + u_{n-1}) + u_n = s_{n-1} + u_n$.
\end{bewijs}
\end{theorie}

\voorbeeld{}{Gegeven: $s_n = n^2 + 2n$. Bepaal $u_n$. Welke soort rij is dat?\label{vb:h3-uitsom}}{%
\uitleg{De eerste term is de eerste som. Voor de andere termen trekken we twee opeenvolgende sommen af. Daarvoor schrijven we $s_{n-1}$ eerst met $n$.}
\begin{stappen}
u_1 = s_1 = 3 & We vervangen $n$ door 1. \\
s_{n-1} = (n-1)^2 + 2(n-1) = n^2 - 1 & We vervangen $n$ door $n - 1$ en werken uit. \\
u_n = n^2 + 2n - (n^2 - 1) = 2n + 1 & Dit geldt voor $n \geqslant 2$. Ook $u_1 = 3$ past in dit voorschrift. \\
\text{rekenkundig met } v = 2 & De rij is $3,\ 5,\ 7,\ \ldots$ Controle: $3 + 5 + 7 = 15 = s_3$.
\end{stappen}}

\begin{oefeningen}
\oef{h3-k7}
\end{oefeningen}

%========================================================================================
\section{Sparen met vaste stortingen}

Wie elk jaar hetzelfde bedrag spaart op een rekening met samengestelde intrest, wil weten hoeveel er na een aantal jaar op staat. Elke storting groeit apart aan, maar niet even lang. Kijk eerst naar een eenvoudig geval: je stort drie jaar lang aan het einde van elk jaar \eur 100, aan 10\,\% intrest. De eerste storting staat er nog twee jaar en wordt $100 \cdot 1{,}1^2 = 121$. De tweede staat er een jaar en wordt $100 \cdot 1{,}1 = 110$. De laatste brengt nog niets op. Samen is dat $100 + 110 + 121 = 331$ euro: de som van een meetkundige rij.

In het algemeen stort je $n$ jaar lang aan het einde van elk jaar een bedrag $a$, aan $p\,\%$ samengestelde intrest, met $q = 1 + \tfrac{p}{100}$:

\vspace{2mm}
\begin{tikzpicture}[x=2.25cm,font=\small]
  \draw[->,lijngrijs] (.55,0) -- (5.6,0) node[right,text=black] {tijd};
  \foreach \x/\t/\w in {1/{jaar 1}/{a\,q^{\,n-1}},2/{jaar 2}/{a\,q^{\,n-2}},4/{jaar $n-1$}/{a\,q},5/{jaar $n$}/{a}}{
    \fill[accent] (\x,0) circle (2pt);
    \node[above=1pt,font=\scriptsize,text=grijs] at (\x,0) {\t};
    \node[below=2pt] at (\x,0) {$\w$};}
  \node[fill=white,inner sep=1pt] at (3,0) {$\cdots$};
\end{tikzpicture}

\begin{theorie}
\begin{stelling}{Eindwaarde van vaste stortingen}
Stort je $n$ jaar lang aan het einde van elk jaar een bedrag $a$ aan $p\,\%$ samengestelde intrest, dan staat er na de laatste storting

\vspace{1mm}
\qquad $\formule{E = a \cdot \dfrac{q^{\,n} - 1}{q - 1}}$ \quad met $\formule{q = 1 + \dfrac{p}{100}}$
\end{stelling}

\begin{bewijs}
Lees de stortingen van de laatste naar de eerste: $E = a + a\,q + \cdots + a\,q^{\,n-1}$. Dat is de som van $n$ termen van een meetkundige rij met eerste term $a$ en quotiënt $q$.
\end{bewijs}
\end{theorie}

\lkr{Dit deel past goed bij Economie-Moderne talen. Het leerplan noemt enkel enkelvoudige en samengestelde intrest als context: sla de opdrachten over sparen over als je dit deel niet behandelt. Leningen komen alleen voor in een opdracht op niveau 3.}

\voorbeeld{}{Je stort 6 jaar lang aan het einde van elk jaar \eur 1000 op een rekening met 3\,\% intrest. Hoeveel staat er na de laatste storting op? Hoeveel daarvan is intrest?\label{vb:h3-sparen}}{%
\uitleg{Dit is de situatie van de formule, met $a = 1000$ en $n = 6$. De intrest is wat er meer op staat dan wat je zelf gestort hebt.}
\begin{stappen}
q = 1 + \dfrac{3}{100} = 1{,}03 & De groeifactor bij 3\,\% intrest. \\[1.5ex]
E = 1000 \cdot \dfrac{1{,}03^6 - 1}{1{,}03 - 1} \approx 6468{,}41 & Na de laatste storting staat er ongeveer \eur 6468,41 op. \\[1.5ex]
6468{,}41 - 6 \cdot 1000 = 468{,}41 & Je stortte zelf \eur 6000, dus \eur 468,41 is intrest.
\end{stappen}}

\nujij{Je stort 10 jaar lang aan het einde van elk jaar \eur 500 op een rekening met 2\,\% intrest. Hoeveel staat er na de laatste storting op?}{$500 \cdot \tfrac{1{,}02^{10} - 1}{0{,}02} \approx 5474{,}86$ euro.}

\begin{oefeningen}
\oef{h3-k8}
\end{oefeningen}

%----------------------------------------------------------------------------------------
% OPDRACHTEN: door elkaar. Niveau 1, dan 2, dan 3.
%----------------------------------------------------------------------------------------
\opdrachtenkop

\lkr{De opdrachten staan bewust door elkaar, met stijgende moeilijkheid. In deze versie staat bij elke opdracht het onderwerp, zodat je kunt kiezen.}

\begin{oefeningen}
\oef{h3-soort}

\oef{h3-rrsom}

\oef{h3-sigma}

\oef{h3-mrsom}

\oef{h3-aantal}

\oef{h3-o1}

\oef{h3-noteer}

\oef{h3-o2}

\oef{h3-o3}

\oef{h3-o4}

\oef{h3-o5}

\oef{h3-o6}

\oef{h3-somindex}

\oef{h3-nuitsom}

\oef{h3-mrcontext}

\oef{h3-uitsom}

\oef{h3-sparen}

\oef{h3-o7}

\oef{h3-o8}

\oef{h3-o9}

\oef{h3-o10}

\oef{h3-o11}

\oef{h3-o12}

\oef{h3-o13}

\oef{h3-o14}

\oef{h3-o15}

\oef{h3-o16}

\oef{h3-o17}

\oef{h3-o18}

\oef{h3-o19}

\oef{h3-o20}

\oef{h3-o21}

\oef{h3-o22}

\oef{h3-o23}
\end{oefeningen}

%----------------------------------------------------------------------------------------
% KOPPEL: verbanden met andere leerstof
%----------------------------------------------------------------------------------------
\opdrachtsectie{Koppel}

\lkr{Kansrekenen, matrices en integralen: kijk na of ze al gezien zijn, of laat die oefeningen weg.}

\begin{oefeningen}
\oef{h3-koppel-1}

\oef{h3-koppel-3}

\oef{h3-koppel-4}
\end{oefeningen}

%----------------------------------------------------------------------------------------
\opdrachtsectie{Test jezelf}

\begin{oefeningen}
\oef{h3-t1}

\oef{h3-t2}

\oef{h3-t3}

\oef{h3-t4}

\oef{h3-t5}

\oef{h3-t6}

\oef{h3-t7}

\oef{h3-t8}

\oef{h3-t9}
\end{oefeningen}

\herhaalschema{%
  \fv{oef:h3-t1}{vb:h3-rr}{oef:h3-rrsom}%
  \fv{oef:h3-t2}{vb:h3-aantal}{oef:h3-aantal}%
  \fv{oef:h3-t3}{vb:h3-mr}{oef:h3-mrsom}%
  \fv{oef:h3-t4}{vb:h3-somteken}{oef:h3-somindex}%
  \fv{oef:h3-t5}{vb:h3-somteken}{oef:h3-noteer}%
  \fv{oef:h3-t6}{vb:h3-nuit}{oef:h3-nuitsom}%
  \fv{oef:h3-t7}{vb:h3-uitsom}{oef:h3-uitsom}%
  \fv{oef:h3-t8}{vb:h3-mr}{oef:h3-mrcontext}%
  \fv{oef:h3-t9}{vb:h3-sparen}{oef:h3-sparen}}

\ifzelfstudie
\section*{Maak je eigen samenvatting}
Sluit dit boek en schrijf in je schrift: de somformule van een rekenkundige en van een meetkundige rij, met het bewijs van elk, hoe je het aantal termen vindt, wat het somteken betekent, hoe je $u_n$ uit $s_n$ haalt, en welke fouten je wilt vermijden. Vergelijk daarna met de kaders.
\fi
